\documentclass{article}
\usepackage[T1]{fontenc}
\usepackage{amsmath}
\usepackage{amssymb}

\title{Assignment 8}

\begin{document}
\maketitle

\section*{Question 1}
 
\subsection*{(a)}

we are given:

\begin{itemize}
    \item pressure: $P=29\,\mathrm{atm}$
    \item temperature: $T=210^\circ\mathrm{C}=483.15\,\mathrm{K}$
    \item heat capacity ratio: $\gamma=1.3$
    \item molar mass of ammonia ($\mathrm{NH_3}$): $M=17.03\times10^{-3}\,\mathrm{kg\,mol^{-1}}$
    \item required resolution: $d=1\,\mathrm{mm}$
\end{itemize}

first we need the speed of sound in the ammonia atmosphere. treating ammonia as an ideal gas,

\[
v=\sqrt{\frac{\gamma P}{\rho}}=\sqrt{\frac{\gamma RT}{M}}
\]

substituting,

\[
v=\sqrt{\frac{(1.3)(8.314)(483.15)}{17.03\times10^{-3}}}
\]
\[
v\approx554\,\mathrm{m\,s^{-1}}
\]

the smallest spot that can be resolved is given by Abbe's limit,

\[
d=\frac{\lambda}{2\sin\theta}
\]

where $\theta$ is the half angle of the diffracted waves collected by the alien's body. the alien is very close to the object ($\ll1\,\mathrm{m}$ while its body is about $1\,\mathrm{m}$ across), so it collects waves from almost all angles up to

\[
\theta\rightarrow90^\circ
\]
\[
\sin\theta\rightarrow1
\]

so,

\[
d=\frac{\lambda}{2}
\]
\[
\lambda=2d=2(1\times10^{-3})=2\times10^{-3}\,\mathrm{m}
\]

using $v=f\lambda$,

\[
f=\frac{v}{\lambda}=\frac{554}{2\times10^{-3}}
\]
\[
f\approx2.77\times10^{5}\,\mathrm{Hz}
\]

\[
\boxed{f\approx277\,\mathrm{kHz}}
\]

so the alien needs to hear up to about 277 kHz.

\subsection*{(b)}

now the object is $2\,\mathrm{m}$ away. the alien's body has diameter $D=1\,\mathrm{m}$, so it only collects waves up to a smaller angle.

\[
\tan\theta=\frac{D}{2f_{\mathrm{d}}}
\]
\[
\theta = \arctan\left(\frac{D}{2f_d}\right)
\]

where $f_{\mathrm{d}}=2\,\mathrm{m}$ is the distance to the object

\[
\theta=\arctan\left(\frac{1}{2\times2}\right)=\arctan(0.25)\approx14.04^\circ
\]
\[
\sin\theta\approx0.2425
\]
 
rearranging Abbe's limit,
 
\[
\lambda=2d\sin\theta
\]
\[
\lambda=2(1\times10^{-3})(0.2425)
\]
\[
\lambda\approx4.85\times10^{-4}\,\mathrm{m}
\]
 
then,
 
\[
f=\frac{v}{\lambda}
=\frac{554}{4.85\times10^{-4}}
\]
\[
f\approx1.14\times10^{6}\,\mathrm{Hz}
\]

therefore, 

\[
\boxed{f\approx1.1\,\mathrm{MHz}}
\]
 
\section*{Question 2}

\subsection*{(a)}

for a diffraction grating at normal incidence,

\[
m\lambda=d\sin\theta
\]

so for each order the diffraction angle is

\[
\theta=\arcsin\left(\frac{m\lambda}{d}\right)
\]
 
we want the angle between $400\,\mathrm{nm}$ and $700\,\mathrm{nm}$ to be as large as possible. using the first order ($m=1$) for the best separation,

\[
\Delta\theta
=
\arcsin\left(\frac{700\,\mathrm{nm}}{d}\right)
-
\arcsin\left(\frac{400\,\mathrm{nm}}{d}\right)
\]

this increases as $d$ gets smaller because the angles get larger. however, we need $\sin\theta\leq1$ for the red light to be diffracted so

\[
\frac{700\,\mathrm{nm}}{d}\leq1
\]
\[
d\geq700\,\mathrm{nm}
\]

the best separation is at the smallest allowed period,

\[
\boxed{d\approx700\,\mathrm{nm}}
\]
 
\subsection*{(b)}
 
using $d=700\,\mathrm{nm}$ and $m=1$:

\begin{itemize}
    \item for the $400\,\mathrm{nm}$ light,

\[
\sin\theta_{400}=\frac{400}{700}=0.5714
\]
\[
\theta_{400}\approx34.85^\circ
\]

\item for the $700\,\mathrm{nm}$ light,

\[
\sin\theta_{700}=\frac{700}{700}=1
\]
\[
\theta_{700}=90^\circ
\]
\end{itemize}

the spectrum is therefore spread over

\[
\Delta\theta=90^\circ-34.85^\circ
\]
\[
\boxed{\Delta\theta\approx55^\circ}
\]

from about $35^\circ$ to $90^\circ$.

\section*{Question 3}

\subsection*{(a)}

the input splits into two waves that go round the loop in opposite directions,

\[
E_A=tE_{in}e^{-i\phi_A}
\qquad
E_B=rE_{in}e^{-i\phi_B}
\]

each output gets one transmitted and one reflected part,
 
\[
E_1=rE_A+tE_B
\qquad
E_2=tE_A-rE_B
\]

the optical path length is the same both ways, so $\phi_A=\phi_B=\phi$,

\[
E_1=2rt\,E_{in}e^{-i\phi}
\qquad
E_2=(t^2-r^2)E_{in}e^{-i\phi}
\]

for $r=t=\frac{1}{\sqrt2}$, $2rt=1$ and $t^2-r^2=0$,

\[
\boxed{E_1=E_{in}, \quad E_2=0}
\]

\subsection*{(b)}

the fibre length is $L=2\pi NR$ and the coupler moves at $\Omega R$. light travels at $c/n$.

let $t_0=\dfrac{2\pi NnR}{c}$ (the travel time with no rotation) and $x=\dfrac{\Omega nR}{c}\ll1$.

\begin{itemize}
    \item clockwise ($t_{CC}$):
    
the coupler moves away from the light, so in time $t_{CC}$ the light must cover the fibre plus the extra distance the coupler moves,

\[
\frac{c}{n}t_{CC}=2\pi NR+\Omega R\,t_{CC}
\]
\[
t_{CC}\left(\frac{c}{n}-\Omega R\right)=2\pi NR
\]
\[
t_{CC}=\frac{2\pi NR}{\frac{c}{n}-\Omega R}
\]

taking out $\frac{c}{n}$ from the bottom,

\[
t_{CC}=\frac{2\pi NnR}{c}\,\frac{1}{1-\frac{\Omega nR}{c}}
=\frac{t_0}{1-x}
\]

since $x\ll1$ we use $\frac{1}{1-x}\approx1+x$,

\[
\boxed{t_{CC}\approx\frac{2\pi NnR}{c}\left(1+\frac{\Omega nR}{c}\right)}
\]

    \item anticlockwise ($t_{AC}$): 
    
the coupler moves towards the light, so the sign of $\Omega$ flips,

\[
\frac{c}{n}t_{AC}=2\pi NR-\Omega R\,t_{AC}
\]
\[
t_{AC}=\frac{2\pi NR}{\frac{c}{n}+\Omega R}
=\frac{t_0}{1+x}
\]

using $\frac{1}{1+x}\approx1-x$,

\[
\boxed{t_{AC}\approx\frac{2\pi NnR}{c}\left(1-\frac{\Omega nR}{c}\right)}
\]

together these are the result we were asked to show,

\[
t=\frac{2\pi NnR}{c}\left(1\pm\frac{\Omega nR}{c}\right)
\]

    \item difference:

\[
\Delta t=t_{CC}-t_{AC}
=t_0(1+x)-t_0(1-x)=2xt_0
\]

\[
\Delta t=\frac{2\pi NnR}{c}\cdot\frac{2\Omega nR}{c}
=\frac{4\pi N\Omega n^2R^2}{c^2}
\]

the area inside the loop is $A=\pi R^2$, so

\[
\boxed{\Delta t=\frac{4\Omega Nn^2A}{c^2}}
\]

\end{itemize}

\subsection*{(c)}

given:

\begin{itemize}
    \item $\lambda = 1550 \text{ nm}$
    \item $n=1.5$
    \item $R=1\text{ m}$
    \item $\Omega=1\,\mathrm{rev/hr}=\frac{2\pi}{3600}=1.745\times10^{-3}\,\mathrm{rad\,s^{-1}}$
    \item $r=t=\frac{1}{\sqrt{2}}$
    \item $A=\pi\,\mathrm{m^2}$,
\end{itemize}
 
\[
\Delta t=\frac{4(1.745\times10^{-3})(1000)(1.5)^2\pi}{(3.00\times10^8)^2}
\approx5.49\times10^{-16}\,\mathrm{s}
\]
 
\[
\omega=\frac{2\pi c}{\lambda}\approx1.22\times10^{15}\,\mathrm{rad\,s^{-1}}
\]
 
\[
\Delta\phi=\omega\Delta t\approx0.667\,\mathrm{rad}
\]
 
for $r=t=\frac{1}{\sqrt2}$ the output is $E_2=\frac{E_{in}}{2}\left(e^{-i\phi_A}-e^{-i\phi_B}\right)$, so
 
\[
\frac{|E_2|^2}{E_{in}^2}=\sin^2\left(\frac{\Delta\phi}{2}\right)
=\sin^2(0.3336)\approx0.107
\]
 
\[
\boxed{\approx11\%}
\]

\end{document}
